\documentclass{article}
    
%Necessario per la traduzione di caratteri in encoding UTF-8 nel linguaggio di LaTeX
    \usepackage[utf8]{inputenc}
    
%Necessario per aggiungere determinati caratteri dell'alfabeto latino e altri simboli come quelli per le frecce
    \usepackage{lmodern,textcomp}

%Necessario per il font utilizzato
    \usepackage{tgheros}
    
%Necessario per la traduzione in italiano
    \usepackage[italian]{babel}

%Necessario per la defizione dei margini e della struttura della pagina
    \usepackage[nomarginpar, left=3cm,right=2cm,top=3cm,bottom=2cm]{geometry}

%Forza posizione immagini
\usepackage{float}

%Necessario per l'utilizzo dello stile di pagina "fancy" che permette l'header e il footer in ogni pagina e mostrare il numero di pagina sul totale
    \usepackage{fancyhdr}
    \usepackage{lastpage}

%Necessari per la creazione dei diagrammi e gli elementi disegnati della pagina
    \usepackage{tikz}
%Aiuta il pacchetto sopra, permettendo l'uso di nodi relativi alla pagina es(\node[anchor=south] at (current page text area.south))
    \usepackage{tikzpagenodes}


%Necessario per permettere di utilizzare il parametro [H] per le tabelle e le figure in modo che vengano posizionate esattamente nel punto indicato 
    \usepackage{float}

%Istruzioni per indicare il numero di livelli di gerarchia per i quali generare la numerazione nel documento(Secnumdepth) e nella table of content (tocdepth)
% Valore 5 : Section -> SubSection -> SubSubSection -> Paragraph -> SubParagraph
    \usepackage{titlesec}
    \setcounter{secnumdepth}{5}
    \setcounter{tocdepth}{5}
%Istruzioni per specificare il desiderio di avere il rientro  e il grassetto anche per i paragrafi e sottoparagrafi
    \titleformat{\paragraph}
        {\normalfont\normalsize\bfseries}{\theparagraph}{1em}{}
    \titlespacing*{\paragraph}{0pt}{3.25ex plus 1ex minus .2ex}{1.5ex plus .2ex}
    \titleformat{\subparagraph}
        {\normalfont\normalsize\bfseries}{\thesubparagraph}{1em}{}
    \titlespacing*{\subparagraph}{0pt}{3.25ex plus 1ex minus .2ex}{.75ex plus .1ex}
    
%Necessario per poter includere hyperlink
    \usepackage[hidelinks]{hyperref}

\pagestyle{fancy}
    \fancyhf{}
    \lhead{A.C. Torre Archimede}
    \rhead{Relazione}
    \rfoot{Pagina \thepage / \pageref{LastPage}}
    
    
\title{
    \includegraphics[scale=0.5]{images/logo.png} \\
    \vspace*{1in}
    {
        \Huge \textbf{A.C. Torre Archimede}}\\
        \vspace*{0.25in}
        \textbf{Relazione Tecnologie Web}\\
        \vspace{0.2in}
        \textbf{Versione: 0.0.2}\\
        \vspace{0.1in}
    }

\author{
    \begin{tabular}[t]{c@{}c}
        \hline
        \\
        Michele Bettin \\
        Marco Marchiante \\
        Stefano Meneguzzo\\
        Edoardo Retis\\
    \end{tabular}
    \vspace*{0.5in} \\
    Dipartimento di Matematica \\
    \textbf{Università degli Studi di Padova} \\
} 
\date{\today}


%--------------------Make usable space all of page
\setlength{\oddsidemargin}{0in} \setlength{\evensidemargin}{0in}
\setlength{\topmargin}{0in}     \setlength{\headsep}{.25in}
\setlength{\textwidth}{6.5in}   \setlength{\textheight}{8.5in}
%--------------------Indention
\setlength{\parindent}{1cm}

\begin{document}
{
    %\fontfamily{qhv}
    \selectfont
    \maketitle
    \newpage
    \tableofcontents
    \newpage

    \section{Abstract}
\subsection{Introduzione}
Il progetto illustrato da questa relazione ha lo scopo di creare un sito internet che rappresenta una fittizia società calcistica, denominata "Associazione Calcistica Torre Archimede". Tale società ha sede in Italia, più precisamente a Padova e opera sotto la giurisdizione della Federazione Italiana Giuoco Calcio (FIGC). La società ha una squadra di calcio che partecipa a campionati italiani organizzati dalla FIGC, uno stadio, denominato "TorreArchimede Stadium", sede della stessa e che usa per allenare la propria squadra, per disputare gare "in casa", e una tifoseria che ogni settimana si riunisce al botteghino per acquistare un biglietto.\\
L'obiettivo del sito è di essere un punto di riferimento sia per la società calcistica sia per i fan che navigano su Internet, cercando di raggiungere il maggior numero di persone. Infatti, il sito è usato per rendere partecipi dell'attività della società ai tifosi che già frequentano lo stadio e, al tempo stesso, permettere ai nuovi tifosi di venire a conoscenza del mondo calcistico del Torre Archimede.\\
Il sito fornisce degli strumenti per permettere a tutte le persone interessate di acquistare online i biglietti di ingresso al campo da gioco, alleggerendo il lavoro di vendita ai tradizionali botteghini.




    \newpage
    \section{Analisi}
    \subsection{Analisi dell'Utenza}
    Il target principale del sito è la tifoseria dell'AC Torre Archimede che, pur essendo abituata al mondo del calcio e familiare con lo stadio della squadra, potrebbe non essere pratica di strumenti informatici. Per questo motivo, il sito internet deve tenerne conto, rendendo il contenuto il più chiaro e semplice possibile.\\
    Un altro punto da considerare sull'utenza del sito è la sua provenienza: la squadra dell'AC Torre Archimede partecipa soltanto a campionati minori, senza mai operare a livello internazionale, presumendo che l'utenza sia italiana, per la maggior parte di loro, residente in zone limitrofe allo stadio. Per questo motivo, la lingua usata nel sito deve essere quella italiana (ad eccezione dei comuni termini di internet, ad esempio, Home, News, Login, ecc.), non richiedendo più versioni alternative in diverse lingue.

    \subsection{Base Informativa}
    Secondo un analisi preliminare del tipo di utenti, le informazioni che un utente si aspetta di trovare e le domande che il sito deve rispondere quando esso vi si accede riguardano l'attività sportiva della società, in particolar modo:
    
    \begin{itemize}
        \item Lo stadio (Dove si trova? Come è fatto?);
        \item La squadra (Chi sono i calciatori? e i dirigenti? A quali eventi partecipano?);
        \item Le gare (Quando giocano? In che stadio? Contro chi? Quanto costa un biglietto?).
    \end{itemize}
    
    \noindent
    Tenendo conto di queste considerazioni, per soddisfare il fabbisogno informativo dell'utente, il sito deve proporre informazioni che riguardano gli argomenti sopra citati, utilizzando come strumenti comunicativi:
    
    \begin{itemize}
        \item Testo descrittivo, se si vuole utilizzare descrizioni per approfondire determinati argomenti (ad esempio, descrivendo come è fatto lo stadio della squadra);
        \item Articoli, se si vuole tenere aggiornato l'utente su fatti che accadono nella sfera quotidiana della società;
        \item Elenchi e tabelle, se si vuole comunicare all'utente una raccolta di dati riferiti ad un certo argomento.
    \end{itemize}


    \subsection{Attori}
    Basandosi sulle necessità del sito, tutti gli utenti si possono raggruppare in tre tipologie: gli utenti non autenticati, gli utenti autenticati e gli amministratori.
    
        \subsubsection{Attore utente non autenticato}
        
        Gli utenti non autenticati sono persone che possono navigare sul sito, ma hanno qualche limitazione. Possono visitare la Homepage, guardare le notizie pubblicate dalla società, leggere la descrizione sullo stadio, consultare l'elenco delle gare in programma, ma non possono acquistare biglietti oppure accedere a sezioni apposite per altri tipi di utenti. Questo tipo di utente può registrarsi al sito per ottenere delle credenziali d'accesso, utili per poter fare l'autenticazione, oppure autenticarsi se ne è già in possesso.
        
        \subsubsection{Attore utente autenticato}
        
        Gli utenti autenticati sono persone che hanno inserito le proprie credenziali d'accesso e si sono autenticati correttamente. Oltre a poter navigare in tutte le aree del sito, possono acquistare biglietti, scegliendo la gara e il settore dove guardarla, e accedere ad una sezione personale dove vengono visualizzati i dati anagrafici dell'utente e l'elenco dei suoi biglietti acquistati.\\ L'utente autenticato può anche togliere l'autenticazione per ritornare ad essere un utente non autenticato.
        
        \subsubsection{Attore amministratore}
        
        Gli amministratori sono utenti autenticati al sito con delle credenziali d'accesso apposite che gli permettono l'ingresso ad una sezione riservata esclusivamente a loro. In questa area, hanno strumenti per:
        \begin{itemize}
            \item Aggiungere articoli inerenti alla società;
            \item Aggiungere o togliere le gare del calendario.
        \end{itemize}
        Anche loro hanno la possibilità di navigare nel sito e fare acquisti di biglietti.
        

    \subsection{Requisiti funzionali}
    Il primo requisito funzionale che il nostro sito deve adempiere è avere una comunicazione efficace delle informazioni da esporre; cioè organizzare il contenuto in modo tale da poter essere facilmente compreso dalla maggior parte degli utenti che navigano in Internet. Per soddisfare questo requisito, si è sviluppato il sito concentrandosi di più sul contenuto da esporre, sviluppando un design minimale, ma efficace e adattabile a qualsiasi dispositivo.\\
    Un altro requisito da soddisfare è mettere a disposizione dell'utente una piattaforma che gli permetta di fare acquisti online per comprare dei biglietti; a questo scopo si è implementato all'interno del sito una serie di strumenti per:
    
    \begin{itemize}
        \item Controllare le partite future e/o da acquistare;
        \item Acquistare un biglietto;
        \item Controllare i biglietti acquistati;
        \item Cancellare e chiedere eventualmente il rimborso per un biglietto acquistato.
    \end{itemize}
    
    \noindent
    L'ultimo requisito da soddisfare, ma non meno importante, è avere un sito facilmente navigabile e accessibile a tutti. Il sito deve essere accessibile per evitare di limitare la navigazione alle persone con difficoltà nel distinguere i colori o alle persone non vedenti.\\
    Per soddisfare questo requisito sono state studiate delle strategie per permettere sia a persone normodotate che diversamente abili di riuscire a trovare facilmente tutte le informazioni e i contenuti all'interno del sito.
    
    


    \newpage
    \section{Organizzazione interna}
    \subsection{Divisione dei compiti}
    Tutti i membri del gruppo hanno avuto modo di toccare tutte le tipologie di lavoro, alternandosi a vicenda nello svolgere i diversi compiti.\\
    Nell'elenco sottostante riportiamo le parti più sostanziali svolte da ogni componente:
    \begin{itemize}
        \item \textbf{Marco}: Back-end PHP, CSS Desktop e Javascript;
        \item \textbf{Michele}: Database, Back-end PHP e CSS Desktop;
        \item \textbf{Edoardo}: Struttura pagine HTML e Relazione;
        \item \textbf{Stefano}: Javascript e CSS Mobile/Tablet/Stampa.
    \end{itemize}
    
    \subsection{Tempistiche}
    Lo sviluppo del progetto è inizialmente cominciato a Dicembre 2021, ma è stato deciso di ricominciarlo da zero a Febbraio 2022, in modo da raggiungere una qualità che potesse soddisfare tutti i membri del gruppo e rispettare tutti i requisiti richiesti dalla consegna.\\
    Da Febbraio fino a Giugno le scadenze sono state rispettate con continuità. Ci sono stati solo dei ritardi nelle ultime settimane, a causa della mancanza di tempo del gruppo, concentrato a studiare per la sessione estiva.
    In conclusione, ci riteniamo soddisfatti del rispetto delle scadenze.

    \newpage
    \section{Tecnologie}
    \subsection{Ambiente di sviluppo}
    Lo sviluppo è stato effettuato su Windows 10/11 e in parte su Ubuntu 21.10 da parte di 4 diversi sviluppatori e con l’utilizzo di Visual Studio Code come IDE.
    Per questa ragione la maggior parte dello sviluppo, del backend è stato eseguito in un ambiente virtuale (un container docker in particolare) con una configurazione il più possibile simile all’ambiente di deploy.
    Ciò unito a uno sviluppo collaborativo tramite Github ha avuto come conseguenza una bassa  incompatibilità con l’ambiente di deploy ee il vantaggio di non dover costantemente fare test di caricamento sul server di laboratorio.
    \subsection{Tecnologie utilizzate}
    Le tecnologie utilizzate per la realizzazione del sito sono state in parte influenzate dall’ambiente finale di deploy (quale il server di laboratorio), in particolare gli applicativi lato Server: PHP (v. 7.4.3), Apache2 (v. 2.4.41), MariaDB (v.10.3.34) e Ubuntu (20.02.4 LTS).\\
    A lato Client è stato utilizzato ovviamente JavaScript e per quanto riguarda HTML e CSS si è optato per HTML5 e CSS3. Nella scelta tra XHTML 4.1 e HTML5 e tra CSS 2 e 3 si è optato per l’ultima versione poichè pur essendo tecnologie nuove e non ancora pienamente supportate, varie fonti (W4C survey e \blue{\url{caniuse.com}}) riportano una diffusione tale da non creare problemi di compatibilità anche con i browser più vecchi a, patto di non utilizzare componenti troppo avanzate dello standard. Oltre alle informazioni reperite in rete abbiamo eseguito anche i nostri piccoli test che hanno confermato le nostre ricerche.\\
    Una volta fatti i conti con la realtà, è risultato che anche i browser più vecchi supportano discretamente lo standard HTML5 (fatta eccezione per IE 6 e 9, che a volte è sono un po’ restrittivi ma in casi non troppo complicati è sempre possibile un workaround), mentre i problemi sono più forti con CSS3 che è ben supportato dalla maggior parte dei browser attualmente in uso, ma ha grosse difficoltà nei browser più vecchi come IE e Opera Mini che supportano molto poco le caratteristiche principali del nuovo standard come flex, grid e attributi avanzati.\\
    Fatte queste considerazioni, è stato deciso di usufruire dei vantaggi di pulizia e raffinatezza offerti dal nuovo standard CSS3.\\
    Per quanto riguarda le versioni del software client (quindi JavaScript e browser), il sito è stato prodotto sulla latest release di Google Chrome ma testato frequentemente su altre versioni spaziando inoltre tra diversi browser (vedi capitolo dei \blue{\hyperref[Test]{Test di accessibilità}}).
    \subsection{Configurazione e installazione}
    Il sito è già configurato e per eseguirlo è sufficiente posizionare i files all'interno della root directory del web server.\\
    Per configurare il database, è sufficiente configurare nel file \texttt{/php/actions/DBAccess.php} i parametri di configurazione e importare il file SQL nel DBMS.\\
    Bisogna poi modificare adeguatamente i files \texttt{.htaccess} e \texttt{/php/helpers/Session.php} (righe 28 e 46), riportando il path assoluto rispetto alla web root del file \texttt{/php/error.php}, per consentire di visualizzare gli errori a prescindere da dalla pagina corrente.
    \newpage
    \section{Progettazione}
    \subsection{Schemi di ordinamento}
    \subsection{Emotional Design}
    \subsection{Accessibilità}
    \subsection{Usabilità}


    \newpage
    \section{Base di Dati}
    \subsection{Scelta delle tabelle}
    \subsection{Identificatori primari}
    \subsection{Forma normale}
    \subsection{Utilizzo atteso}

    
    \newpage
    \section{Costruzione}
    \subsection{Separazione}
    \subsection{Presentazione}
    \label{Presentazione}
        \subsubsection{CSS Desktop}
        \subsubsection{CSS Mobile}
        \subsubsection{CSS di Stampa}
    \subsection{Struttura}
    \subsection{Comportamento}


    \newpage
    \section{Strutturazione}
    \subsection{Sito}
    \subsection{Pagine}


    \newpage
    \input{9_Accessibilità}
    \newpage
    \input{10_Usabilità}
    \newpage
    \section{Conclusioni}


    \newpage

    

}

\end{document}